\documentclass[aspectratio=169, 10pt]{beamer}

%----------------------------------------------------------------------------------------
%   PACKAGES
%----------------------------------------------------------------------------------------
\usepackage[utf8]{inputenc}
\usepackage[T1]{fontenc}
\usepackage{lmodern}
\usepackage{amsmath, amssymb, amsthm}
\usepackage{tikz}

%----------------------------------------------------------------------------------------
%   COLOR DEFINITIONS
%----------------------------------------------------------------------------------------
% Standard Transgender Pride Flag Hex Codes
\definecolor{TransBlue}{HTML}{5BCEFA}
\definecolor{TransPink}{HTML}{F5A9B8}
\definecolor{TransWhite}{HTML}{FFFFFF}
\definecolor{TransBlue}{HTML}{5BCEFA}
\definecolor{TransPink}{HTML}{F5A9B8}
\definecolor{TransWhite}{HTML}{FFFFFF}
\definecolor{tyellow}{HTML}{FFCC5B}

% Auxiliary Colors for Readability and Depth
\definecolor{DarkText}{HTML}{2B2B2B}      % Dark Charcoal for main text
\definecolor{DeepBlue}{HTML}{1A7F9B}      % Darker version of TransBlue for structural text
\definecolor{DeepPink}{HTML}{B05E6B}      % Darker version of TransPink for alerts
\definecolor{LightGrey}{HTML}{F0F0F0}     % For block backgrounds

%----------------------------------------------------------------------------------------
%   THEME SETTINGS
%----------------------------------------------------------------------------------------

% Disable navigation symbols for a cleaner look
\setbeamertemplate{navigation symbols}{}

% font theme
\usefonttheme{professionalfonts} % Best for math
\setbeamerfont{title}{series=\bfseries, size=\Large}
\setbeamerfont{frametitle}{series=\bfseries}

% 1. Palette Colors
% Primary: Used for title background, section headers
\setbeamercolor{palette primary}{bg=TransBlue, fg=DarkText}
% Secondary: Used for subsections
\setbeamercolor{palette secondary}{bg=TransPink, fg=DarkText}
% Tertiary: Used for outer theme elements
\setbeamercolor{palette tertiary}{bg=TransWhite, fg=DarkText}
% Quaternary: Title bar in standard themes
\setbeamercolor{palette quaternary}{bg=DarkText, fg=TransWhite}

% 2. Structure (Itemize bullets, Table of Contents)
\setbeamercolor{structure}{fg=DeepBlue}
\setbeamercolor{local structure}{fg=DeepBlue}

% 3. Standard Text
\setbeamercolor{normal text}{fg=DarkText, bg=white}

% 4. Block Environments (Theorems, Definitions, Proofs)
% Standard Block (Theorems)
\setbeamercolor{block title}{bg=TransBlue, fg=DarkText}
\setbeamercolor{block body}{bg=TransBlue!15!white, fg=DarkText}

% Alert Block (Important/Warning)
\setbeamercolor{block title alert}{bg=TransPink, fg=DarkText}
\setbeamercolor{block body alert}{bg=TransPink!15!white, fg=DarkText}

% Example Block
\setbeamercolor{block title example}{bg=LightGrey, fg=DeepBlue}
\setbeamercolor{block body example}{bg=LightGrey!40!white, fg=DarkText}

% 5. Title Page Customization
\setbeamertemplate{title page}{
    \begin{center}
        \vspace{1cm}
        {\usebeamerfont{title}\usebeamercolor[fg]{structure}\inserttitle}\par
        \vspace{0.5cm}
        {\usebeamerfont{subtitle}\insertsubtitle}\par
        \vspace{1cm}
        {\usebeamerfont{author}\insertauthor}\par
        \vspace{0.2cm}
        {\usebeamerfont{institute}\insertinstitute}\par
        \vspace{0.5cm}
        {\usebeamerfont{date}\insertdate}\par
    \end{center}
    % Decorative Flag Stripes at the bottom
    \begin{tikzpicture}[remember picture, overlay]
        \node[anchor=south west, inner sep=0] at (current page.south west) {
            \begin{tikzpicture}
                \fill[TransBlue] (0,0) rectangle (\paperwidth, 0.25cm);
                \fill[TransPink] (0,0.25) rectangle (\paperwidth, 0.50cm);
                \fill[TransWhite] (0,0.50) rectangle (\paperwidth, 0.75cm);
                \fill[TransPink] (0,0.75) rectangle (\paperwidth, 1.00cm);
                \fill[TransBlue] (0,1.00) rectangle (\paperwidth, 1.25cm);
            \end{tikzpicture}
        };
    \end{tikzpicture}
}

% 6. Frame Title Customization
\setbeamertemplate{frametitle}{
    \nointerlineskip
    \begin{beamercolorbox}[wd=\paperwidth,ht=2.5ex,dp=1.125ex,leftskip=0.5cm]{palette primary}
        \usebeamerfont{frametitle}\insertframetitle
    \end{beamercolorbox}
    \nointerlineskip
    % Thin pink line under the blue header
    \begin{beamercolorbox}[wd=\paperwidth,ht=0.5ex]{palette secondary}
    \end{beamercolorbox}
}

% 7. Footline Customization
\setbeamertemplate{footline}{
    \leavevmode%
    \hbox{%
    \begin{beamercolorbox}[wd=.333\paperwidth,ht=2.25ex,dp=1ex,center]{palette primary}%
        \insertsection
    \end{beamercolorbox}%
    \begin{beamercolorbox}[wd=.333\paperwidth,ht=2.25ex,dp=1ex,center]{palette secondary}%
        \insertsubsection
    \end{beamercolorbox}%
    \begin{beamercolorbox}[wd=.333\paperwidth,ht=2.25ex,dp=1ex,rightskip=2ex]{palette tertiary}%
        \hfill \insertframenumber{} / \inserttotalframenumber
    \end{beamercolorbox}}%
    \vskip0pt%
}


%----------------------------------------------------------------------------------------
%   DOCUMENT CONTENT
%----------------------------------------------------------------------------------------

\title{Transgender Pride Theme}
\subtitle{A Formal Beamer Theme for Mathematical Presentations}
\author{Research Assistant}
\institute{Department of Mathematics}
\date{\today}

\begin{document}

% Slide 1: Title
\begin{frame}[plain]
    \titlepage
\end{frame}

% Slide 2: Table of Contents
\begin{frame}{Outline}
    \tableofcontents
\end{frame}

\section{Introduction}

% Slide 3: Lists and Text
\begin{frame}{Motivation and Palette}
    This theme utilizes the specific hex codes associated with the transgender pride flag.
    \begin{itemize}
        \item \textcolor{DeepBlue}{\textbf{Blue (5BCEFA)}}: Used for primary headers and standard blocks.
        \item \textcolor{DeepPink}{\textbf{Pink (F5A9B8)}}: Used for secondary elements and alert blocks.
        \item \textcolor{DarkText}{\textbf{White (FFFFFF)}}: Used for background and negative space.
    \end{itemize}
    
    \vspace{0.5cm}
    It is important that mathematical content remains legible:
    \begin{enumerate}
        \item Text color is set to dark charcoal (\texttt{\#2B2B2B}) rather than black for softness.
        \item Contrast is maintained on pastel backgrounds.
    \end{enumerate}
\end{frame}

\section{Mathematical Examples}

% Slide 4: Block Environments (Theorems)
\begin{frame}{Standard Environments}
    \begin{definition}[Topological Space]
        A topological space is a pair $(X, \tau)$, where $X$ is a set and $\tau$ is a collection of subsets of $X$, satisfying specific axioms regarding unions and intersections.
    \end{definition}

    \begin{theorem}[Heine-Borel]
        For a subset $S \subseteq \mathbb{R}^n$, the following are equivalent:
        \begin{enumerate}
            \item $S$ is closed and bounded.
            \item $S$ is compact.
        \end{enumerate}
    \end{theorem}
\end{frame}

% Slide 5: Proof Environment
\begin{frame}{Proof and Alerts}
    \begin{alertblock}{Important Note}
        Care must be taken when extending this theorem to infinite-dimensional spaces.
    \end{alertblock}

    \begin{proof}
        Let $\{U_\alpha\}_{\alpha \in A}$ be an open cover of $S$. Since $S$ is bounded, it is contained within a closed hypercube $H$. By the Tychonoff product theorem... \textit{(proof continues)}.
        \qed
    \end{proof}
\end{frame}

% Slide 6: Mathematics Mode
\begin{frame}{Mathematical Typography}
    The theme supports standard equations clearly.
    
    \begin{equation}
        \int_{-\infty}^{\infty} e^{-x^2} \, dx = \sqrt{\pi}
    \end{equation}
    
    Let $f: A \to B$ be a homomorphism between two algebraic structures.
    \[
        \ker(f) = \{ a \in A \mid f(a) = e_B \}
    \]
\end{frame}

\end{document}
